\section{Первое задание}
Для указанной функции провести модульное тестирование разложения
функции в степенной ряд.
Выбрать достаточное тестовое покрытие.
Функция \( \sin x \).

Разложим функцию в ряд Тейлора в окрестности точки \(x = 0 \).

\[
	\sin x =
	\sum_{n = 0}^\infty \frac {(-1)^n}{2n+1)!} \cdot x^{2n + 1} =
	x - \frac{x^3}{3!} + \frac{x^5}{5!} - \frac{x^7}{7!} + \frac{x^9}{9!} - \cdots
\]

С помощью тривиальной математики реализуем вычисление синуса с заданной точностью:

\begin{minted}{python}
def sin_power_series(x, eps=1e-10):
    """
    Вычисляет sin(x) по степенному ряду:
    Итерации продолжаются, пока абсолютное значение последнего слагаемого > eps.
    По очевидным причинам этого достаточно, чтобы полученное значение
    не отходило от настоящего значения синуса больше, чем на eps.
    """
    term = x
    result = x
    n = 1
    while abs(term) > eps:
        term = -term * x * x / ((2 * n) * (2 * n + 1))
        result += term
        n += 1
    return result
\end{minted}

Для данной функции можно предложить следующее покрытие.
Точность \texttt{eps} в данном случае можно зафиксировать произвольную.
\begin{itemize}
	\item Тест при известных из курса тригонометрии
	      значениях \(x\):
	      \(x \in \{ 0, \pi/6, \pi/4, \pi/3, \pi/2, \ldots \} \)
	\item Так как функция \( \sin x \) --- нечетная, можно применить
	      метод property-based тестирования и проверить это свойство
	      для произвольных \(x \).
	      \[ \sin (-x) = - \sin x \]
	\item Также можно воспользоваться свойством периодичности функции:
	      \[ \sin x = \sin (x + 2 \pi) \]
	\item Если мы доверяем реализации функции \( \sin \) в стандартной библиотеки
	      Python, то можно проверить их равенство с помощью property-based testing.
	      \[ | \sin_\varepsilon^* x - \sin x | \leq \varepsilon \]
	      где \( \sin_\varepsilon^* \) -- результат вызова функции \texttt{sin\_power\_series} с соответствующими аргументами
\end{itemize}
